% ==============================================================================
% ОГЛАВЛЕНИЕ ИСХОДНИКА (Word TOC)
% Партия: B03-002 (DISS-005-B0033 .. DISS-005-B0113, 81 блок)
% Точный перенос структуры и нумерации страниц исходника v0.5
% ==============================================================================

% DISS-005-B0033
\chapter*{Содержание\commentnote{DISS-005-C0002}{Алексей Евдаков}{Уровень вложенности и прочие аспекты можно будет по итогу поменять.}}
\markboth{Содержание}{Содержание}
\vspace{0.5cm}

% DISS-005-B0034
\noindent\textbf{Введение}\dotfill\textbf{6}\par\vspace{1.5mm}
% DISS-005-B0035
\noindent\textbf{Глава 1 СОВРЕМЕННЫЕ ПРОБЛЕМЫ БЫСТРОДЕЙСТВУЮЕГО АВТОМАТИЧЕСКОГО ВВОДА РЕЗЕРВА И ВАРИАНТЫ ИХ УСТРАНИЕНИЯ СТАНДАРТНЫМИ МЕТОДАМИ}\dotfill\textbf{10}\par\vspace{1.5mm}
% DISS-005-B0036
\noindent\hspace*{1.5em}1.1. Введение\dotfill 10\par\vspace{0.5mm}
% DISS-005-B0037
\noindent\hspace*{1.5em}1.2. Статистика функционирования БАВР\dotfill 13\par\vspace{0.5mm}
% DISS-005-B0038
\noindent\hspace*{1.5em}1.3. Базовые проблемы функционирования РНМ\dotfill 14\par\vspace{0.5mm}
% DISS-005-B0039
\noindent\hspace*{3.0em}1.3.1. Проблемы базовых органов реле направления мощности\dotfill 14\par\vspace{0.5mm}
% DISS-005-B0040
\noindent\hspace*{3.0em}1.3.2. Внедрение адаптивных функций\dotfill 18\par\vspace{0.5mm}
% DISS-005-B0041
\noindent\hspace*{1.5em}1.4. Анализ органов РНМ других производителей\dotfill 20\par\vspace{0.5mm}
% DISS-005-B0042
\noindent\hspace*{3.0em}1.4.1. ЭКРА\dotfill 20\par\vspace{0.5mm}
% DISS-005-B0043
\noindent\hspace*{3.0em}1.4.2. Механотроника (БМРЗ)\dotfill 21\par\vspace{0.5mm}
% DISS-005-B0044
\noindent\hspace*{3.0em}1.4.3. Бреслер\dotfill 23\par\vspace{0.5mm}
% DISS-005-B0045
\noindent\hspace*{3.0em}1.4.4. ЧЭАЗ\dotfill 24\par\vspace{0.5mm}
% DISS-005-B0046
\noindent\hspace*{3.0em}1.4.5. Радиус Автоматика (Сириус)\dotfill 24\par\vspace{0.5mm}
% DISS-005-B0047
\noindent\hspace*{3.0em}1.4.6. Релематика (ТОР300)\dotfill 24\par\vspace{0.5mm}
% DISS-005-B0048
\noindent\hspace*{3.0em}1.4.7. ABB (SUE 3000)\dotfill 24\par\vspace{0.5mm}
% DISS-005-B0049
\noindent\hspace*{3.0em}1.4.8. Digital Motor Bus Transfer System (М-4172 и М-4272)\dotfill 25\par\vspace{0.5mm}
% DISS-005-B0050
\noindent\hspace*{1.5em}1.5. Сложные случаи функционирования адаптивного РНМ и направления развития\dotfill 26\par\vspace{0.5mm}
% DISS-005-B0051
\noindent\hspace*{3.0em}1.5.1. Адаптивность по времени РНМ (колебание)\dotfill 26\par\vspace{0.5mm}
% DISS-005-B0052
\noindent\hspace*{3.0em}1.5.2. Отпадание адаптивного РНМ (вопрос подбора К отс)\dotfill 28\par\vspace{0.5mm}
% DISS-005-B0053
\noindent\hspace*{3.0em}1.5.3. Блокировка работы БАВР по РНМ при внешних двухфазных замыканиях\dotfill 31\par\vspace{0.5mm}
% DISS-005-B0054
\noindent\hspace*{3.0em}1.5.4. Измерение тока с применением катушек Роговского и ёмкостных делителей\dotfill 39\par\vspace{0.5mm}
% DISS-005-B0055
\noindent\hspace*{3.0em}1.5.5. Необходимость проработки РНМ по 0-ой последовательности для БАВР\dotfill 43\par\vspace{0.5mm}
% DISS-005-B0056
\noindent\hspace*{1.5em}1.6. Выводы\dotfill 44\par\vspace{0.5mm}
% DISS-005-B0057
\noindent\textbf{Глава 2 СОВРЕМЕННЫЕ МОДЕРНИЗАЦИИ БЫСТРОДЕЙСТВУЮЩЕГО АВТОМАТИЧЕСКОГО ВВОДА РЕЗЕРВА СТАНДАРТНЫМИ МЕТОДАМИ}\dotfill\textbf{46}\par\vspace{1.5mm}
% DISS-005-B0058
\noindent\hspace*{1.5em}2.1. Введение\dotfill 46\par\vspace{0.5mm}
% DISS-005-B0059
\noindent\hspace*{1.5em}2.2. Функционирования РНМ с катушками Роговского\dotfill 47\par\vspace{0.5mm}
% DISS-005-B0060
\noindent\hspace*{3.0em}2.2.1. Теоретическое исследование и сравнение КР и ТТ\dotfill 47\par\vspace{0.5mm}
% DISS-005-B0061
\noindent\hspace*{3.0em}2.2.2. Оценка переходного процесса и влияния различных факторов\dotfill 51\par\vspace{0.5mm}
% DISS-005-B0062
\noindent\hspace*{3.0em}2.2.3. Анализ применения преобразования Фурье\dotfill 55\par\vspace{0.5mm}
% DISS-005-B0063
\noindent\hspace*{3.0em}2.2.4. Анализ работы реле направления мощности с применением ТТ и КР\dotfill 57\par\vspace{0.5mm}
% DISS-005-B0064
\noindent\hspace*{3.0em}2.2.5. Практическое применение в промышленности\dotfill 62\par\vspace{0.5mm}
% DISS-005-B0065
\noindent\hspace*{3.0em}2.2.6. Выводы о применении катушек Роговского для БАВР\dotfill 64\par\vspace{0.5mm}
% DISS-005-B0066
\noindent\hspace*{1.5em}2.3. Совершенствование алгоритмов функционирования восстановления нормального режима\dotfill 66\par\vspace{0.5mm}
% DISS-005-B0067
\noindent\hspace*{3.0em}2.3.1. Существующие требования к ВНР\dotfill 67\par\vspace{0.5mm}
% DISS-005-B0068
\noindent\hspace*{3.0em}2.3.2. Последовательный и параллельный ВНР\dotfill 67\par\vspace{0.5mm}
% DISS-005-B0069
\noindent\hspace*{3.0em}2.3.3. Блокировка по однофазным замыканиям на землю\dotfill 67\par\vspace{0.5mm}
% DISS-005-B0070
\noindent\hspace*{3.0em}2.3.4. Блокировка при превышении угла\dotfill 68\par\vspace{0.5mm}
% DISS-005-B0071
\noindent\hspace*{3.0em}2.3.5. Блокировка по частоте\dotfill 69\par\vspace{0.5mm}
% DISS-005-B0072
\noindent\hspace*{3.0em}2.3.6. Уменьшение одновременного срабатывания ВНР\dotfill 70\par\vspace{0.5mm}
% DISS-005-B0073
\noindent\hspace*{3.0em}2.3.7. Выводы о совершенствование алгоритмов функционирования восстановления нормального режима\dotfill 71\par\vspace{0.5mm}
% DISS-005-B0074
\noindent\hspace*{1.5em}2.4. Выводы\dotfill 72\par\vspace{0.5mm}
% DISS-005-B0075
\noindent\textbf{Глава 3 ПОДГОТОВКА ДАТАСЕТА НА ОСНОВЕ РЕАЛЬНЫХ ОСЦИЛЛОГРАММ}\dotfill\textbf{74}\par\vspace{1.5mm}
% DISS-005-B0076
\noindent\hspace*{1.5em}3.1. Введение\dotfill 74\par\vspace{0.5mm}
% DISS-005-B0077
\noindent\hspace*{1.5em}3.2. Методология формирования датасета реальных осциллограмм\dotfill 77\par\vspace{0.5mm}
% DISS-005-B0078
\noindent\hspace*{3.0em}3.2.1. Этапы подготовки датасета\dotfill 77\par\vspace{0.5mm}
% DISS-005-B0079
\noindent\hspace*{3.0em}3.2.2. Проблемы и вызовы при работе с реальными осциллограммами\dotfill 78\par\vspace{0.5mm}
% DISS-005-B0080
\noindent\hspace*{1.5em}3.3. Сбор и первичная обработка осциллограмм\dotfill 80\par\vspace{0.5mm}
% DISS-005-B0081
\noindent\hspace*{3.0em}3.3.1. Источники данных и типы осциллограмм\dotfill 80\par\vspace{0.5mm}
% DISS-005-B0082
\noindent\hspace*{3.0em}3.3.2. Автоматизированный поиск и сбор осциллограмм\dotfill 81\par\vspace{0.5mm}
% DISS-005-B0083
\noindent\hspace*{3.0em}3.3.3. Формирование единой базы данных осциллограмм\dotfill 82\par\vspace{0.5mm}
% DISS-005-B0084
\noindent\hspace*{3.0em}3.3.4. Поиск осциллограмм, относящихся к конкретным терминалам/объектам\dotfill 85\par\vspace{0.5mm}
% DISS-005-B0085
\noindent\hspace*{1.5em}3.4. Предобработка осциллограмм: деперсонализация и унификация\dotfill 86\par\vspace{0.5mm}
% DISS-005-B0086
\noindent\hspace*{3.0em}3.4.1. Ключевые аспекты формата COMTRADE для дальнейшей обработки\dotfill 86\par\vspace{0.5mm}
% DISS-005-B0087
\noindent\hspace*{3.0em}3.4.2. Необходимость деперсонализации данных\dotfill 87\par\vspace{0.5mm}
% DISS-005-B0088
\noindent\hspace*{3.0em}3.4.3. Процесс деперсонализации\dotfill 87\par\vspace{0.5mm}
% DISS-005-B0089
\noindent\hspace*{3.0em}3.4.4. Унификация параметров осциллограмм\dotfill 88\par\vspace{0.5mm}
% DISS-005-B0090
\noindent\hspace*{1.5em}3.5. Нормализация сигналов и преобразование форматов\dotfill 96\par\vspace{0.5mm}
% DISS-005-B0091
\noindent\hspace*{3.0em}3.5.1. Методика определения базисных значений для нормализации\dotfill 96\par\vspace{0.5mm}
% DISS-005-B0092
\noindent\hspace*{3.0em}3.5.2. Применение нормализации и преобразование осциллограмм в формат CSV\dotfill 99\par\vspace{0.5mm}
% DISS-005-B0093
\noindent\hspace*{1.5em}3.6. Разметка (аннотирование) осциллограмм\dotfill 101\par\vspace{0.5mm}
% DISS-005-B0094
\noindent\hspace*{3.0em}3.6.1. Необходимость разметки для обучения ML моделей и анализа\dotfill 102\par\vspace{0.5mm}
% DISS-005-B0095
\noindent\hspace*{3.0em}3.6.2. Подготовка данных для ручной разметки\dotfill 103\par\vspace{0.5mm}
% DISS-005-B0096
\noindent\hspace*{3.0em}3.6.3. Автоматизация процесса разметки и перспективы\dotfill 103\par\vspace{0.5mm}
% DISS-005-B0097
\noindent\hspace*{3.0em}3.6.4. Фильтрация «пустых» осциллограмм\dotfill 103\par\vspace{0.5mm}
% DISS-005-B0098
\noindent\hspace*{1.5em}3.7. Выводы\dotfill 106\par\vspace{0.5mm}
% DISS-005-B0099
\noindent\textbf{Глава 4 Статистический анализ функционирования БАВР на основе реальных осциллограмм (Планируемая)}\dotfill\textbf{107}\par\vspace{1.5mm}
% DISS-005-B0100
\noindent\hspace*{1.5em}4.1. Введение\dotfill 107\par\vspace{0.5mm}
% DISS-005-B0101
\noindent\textbf{Глава 5 Разработка и исследование адаптивных алгоритмов БАВР с применением методов машинного обучения (Планируемая)}\dotfill\textbf{108}\par\vspace{1.5mm}
% DISS-005-B0102
\noindent\hspace*{1.5em}5.1. Введение\dotfill 108\par\vspace{0.5mm}
% DISS-005-B0103
\noindent\textbf{БИБЛИОГРАФИЧЕСКИЙ СПИСОК}\dotfill\textbf{109}\par\vspace{1.5mm}
% DISS-005-B0104
\noindent\textbf{Приложение 1. Что-то буду сюда добавлять постепенно.}\dotfill\textbf{116}\par\vspace{1.5mm}
% DISS-005-B0105
\noindent\textbf{Глава 6 Мысли в слух}\dotfill\textbf{117}\par\vspace{1.5mm}
% DISS-005-B0106
\noindent\hspace*{1.5em}6.1. ИИ\dotfill 117\par\vspace{0.5mm}
% DISS-005-B0107
\noindent\hspace*{3.0em}6.1.1. Общие знания про ML\dotfill 117\par\vspace{0.5mm}
% DISS-005-B0108
\noindent\hspace*{3.0em}6.1.2. Подготовка датасета\dotfill 118\par\vspace{0.5mm}
% DISS-005-B0109
\noindent\hspace*{3.0em}6.1.3. Перенос идей из LLM в текущие задачи\dotfill 118\par\vspace{0.5mm}
% DISS-005-B0110
\noindent\hspace*{3.0em}6.1.4. Дикие идеи по ИИ\dotfill 121\par\vspace{0.5mm}
% DISS-005-B0111
\noindent\hspace*{1.5em}6.2. БАВР\dotfill 122\par\vspace{0.5mm}
% DISS-005-B0112
\noindent\hspace*{3.0em}6.2.1. Задачи\dotfill 122\par\vspace{0.5mm}

% DISS-005-B0113 (пустая строка завершения оглавления)
\vspace{0.5cm}
